\section{Crystalline realization and the section problem (Part J)}\label{sec:categoryJ}

\subsection{Data, conventions, and limits of the supplied witnesses}
We use $K$ for a complete discretely valued field of characteristic zero with
perfect residue field $k$ of characteristic $p>0$, $W=W(k)$, and $L=W[1/p]$.
We assume that $K/L$ is finite and totally ramified.  Put
$\pi_K=\operatorname{Gal}(\overline K/K)$ and $C=\widehat{\overline K}$.
These are the fields on cote 161-5, scan 023.  In particular, that page
assumes that $k$ is perfect, not that it is algebraically closed.
We normalize $v_p(p)=1$.

There is a serious witness limitation.  The supplied archive contains the
Pro transcription of 161-5/023--029, but not the second transcription
\texttt{K161-5-02} cited in the study, and it contains no images of these
pages.  The following is a typed reconstruction of the mathematics, not a
certification of every uncertain manuscript symbol.  On scan 023 the symbol
$\mathcal M_K^{(1)}$ denotes all Hodge--Tate representations with weights
$0,1$, while $\widetilde{\mathrm{BT}}_V$ is only a full subcategory of it.
The subsequent definition says that $\mathcal M_K^{\mathrm{BT}}$ is generated
by $\mathcal M_K^{(1)}$.  This is different from being generated by actual
Barsotti--Tate representations.  We keep that distinction explicit.

Let $\mathcal C$ be a neutral $\mathbf Q_p$-linear Tannakian category of
representations, with forgetful functor $\omega_p$ and
$G=\operatorname{Aut}^{\otimes}(\omega_p)$, a pro-algebraic group over
$\mathbf Q_p$.  Evaluation gives
$\rho:\pi_K\longrightarrow G(\mathbf Q_p)$.
On $C$ there are two fibre functors: $C\otimes\omega_p$ with its scalar
semilinear action, and the same vector spaces with the diagonal action
induced by $\rho$.  Their tensor-isomorphism torsor is the right $G_C$-torsor
$\Phi_C$.  Its inner automorphism group is
$G_C^!=\operatorname{Aut}_{G_C}(\Phi_C)$.
For Hodge--Tate objects the weight decomposition defines a cocharacter
$j^!:\mathbf G_{m,C}\to G_C^!$.  If $\tau_C$ is the torsor of bases of the
Tate line, with the dual convention chosen to cancel the Hodge--Tate weights,
then
\[
 Q_C=\tau_C\mathbin{\times^{\mathbf G_{m,C}}}\Phi_C,
 \qquad G''_C=\operatorname{Aut}_{G_C}(Q_C)
\]
encode the graded Hodge--Tate fibre functor.  The twist gives a cocharacter
$j_C:\mathbf G_{m,C}\to G''_C$.
The descent of $(Q_C,j_C)$ to a right $G_L$-torsor $Q_L$ and a cocharacter
$j_L$ is an additional datum on scans 025--026, labelled an axiom there;
it does not follow merely from the notation.

Let $U''_L$ be the unipotent subgroup of $G''_L$ acting trivially on the
associated graded for the filtration defined by $j_L$.  A right
$U''_L$-torsor $R_L$ produces
\[
 P_L=R_L\mathbin{\times^{U''_L}}Q_L,
 \qquad G'_L=\operatorname{Aut}_{G_L}(P_L).
\]
The associated filtered fibre functor is
$V\mapsto P_L\times^{G_L}(L\otimes_{\mathbf Q_p}\omega_p(V))$,
with the filtration transported from $Q_L$.  For a crystalline realization
the filtration naturally lives after extension of scalars to $K$, while
Frobenius is an isomorphism $\varphi^*P_L\simeq P_L$ over $L$, where
$\varphi$ is Witt Frobenius.  Scan 027 instead reads ``descent from $L$ to
$K$'' and an isomorphism with differently subscripted torsors.  Since
$L\subset K$, that sentence does not type-check as descent under the
notation of scan 023.  Moreover scan 026 has changed $\pi$ to
$\operatorname{Aut}_L(C)$.  These changes cannot be repaired uniquely from
the single supplied witness.  The section theorem below applies to every
consistent realization of the resulting extension and does not depend on
choosing one of these uncertain readings.

\subsection{The complete abstract classification}
Let $S$ be a ring with an action of a group $\Gamma$, let $G_S$ be a group
scheme with compatible semilinear $\Gamma$-action, and let $P$ be a right
$G_S$-torsor.  Assume the semilinear lifts exist for each element of
$\Gamma$.  The group $E$ of such lifts fits into
\[
 1\longrightarrow H\longrightarrow E\longrightarrow\Gamma
 \longrightarrow1,\qquad H=\operatorname{Aut}_{G_S}(P)(S).
\]
This is the extension of cote 161-5/028.  For topological groups all sections
and cocycles below are required to be continuous; the same proofs apply.

\begin{theorem}[Proved: exact classification of extensions]\label{cat:J1}
Fix a subgroup $\pi\subset\Gamma$ and a section
$\alpha_0:\pi\to E|_\pi$.  If there is no extending section, the solution
set is empty.  If an extending section $\alpha:\Gamma\to E$ exists, let
$\gamma\cdot h=\alpha(\gamma)h\alpha(\gamma)^{-1}$ be the resulting action
on $H$.  Then the groupoid of extending sections, with conjugations that
preserve $\alpha_0$, is
\[
 \bigl[Z^1(\Gamma,H;\pi)/H^\pi\bigr],
 \quad
 Z^1(\Gamma,H;\pi)=
 \{c:\Gamma\to H:\ c(\gamma\delta)=c(\gamma)
       (\gamma\cdot c(\delta)),\ c|_\pi=1\}.
\]
Here $h\in H^\pi$ acts by
$c(\gamma)\mapsto h c(\gamma)(\gamma\cdot h^{-1})$.
If $\pi$ is normal, this groupoid is canonically
$[Z^1(\Gamma/\pi,H^\pi)/H^\pi]$; its set of isomorphism classes is
$H^1(\Gamma/\pi,H^\pi)$.  The identifications depend on the reference
section $\alpha$; changing it twists the coefficient action.
\end{theorem}
\begin{proof}
For another section $\beta$ put
$c(\gamma)=\beta(\gamma)\alpha(\gamma)^{-1}\in H$.
Multiplying $\beta(\gamma)\beta(\delta)$ gives the displayed cocycle
identity.  Conversely that identity says precisely that
$\beta(\gamma)=c(\gamma)\alpha(\gamma)$ is a homomorphism section.
Equality with $\alpha_0$ on $\pi$ is precisely $c|_\pi=1$.
Conjugating $\beta$ by $h$ gives the stated formula.  Such a conjugation
preserves $\alpha_0$ if and only if $h$ commutes with $\alpha_0(\pi)$,
which is the condition $h\in H^\pi$.  Thus even a conjugation allowed
initially in all of $H$ lies in $H^\pi$ when its source and target both
extend the fixed $\alpha_0$.

Assume $\pi$ is normal.  For $n\in\pi$, the cocycle identity gives
$c(\gamma n)=c(\gamma)$ and
$c(n\gamma)=n\cdot c(\gamma)$.  Since
$n\gamma=\gamma(\gamma^{-1}n\gamma)$, the first identity shows that
$n\cdot c(\gamma)=c(\gamma)$.  Consequently $c$ takes values in $H^\pi$
and is constant on cosets.  The $\Gamma$-action on $H^\pi$ factors
through $\Gamma/\pi$, so $c$ descends to a cocycle on the quotient.
Inflation is inverse to this construction and preserves the action of
$H^\pi$.  For continuous cocycles the descended map is continuous by
the quotient topology.  This proves all assertions, including morphisms
and automorphisms in the groupoids.
\end{proof}

\begin{corollary}[Proved: Frobenius classes]\label{cat:J2}
Suppose $\pi\triangleleft\Gamma$, $\Gamma/\pi=\mathbf Z$, and
$H^\pi=G'_L(L)$.  Write $\sigma$ for the action of $1\in\mathbf Z$ on
$G'_L(L)$ induced by a fixed extending section.  The extensions are
classified by the $\sigma$-conjugacy classes
\[
 G'_L(L)/\bigl(b\sim h b\sigma(h)^{-1}\bigr).
\]
The automorphism group of the class represented by $b$ is
$\{h\in G'_L(L):h b=b\sigma(h)\}$.
\end{corollary}
\begin{proof}
A cocycle on the infinite cyclic group is uniquely determined by its
value $b$ at $1$; at $n>0$ it is
$b\sigma(b)\cdots\sigma^{n-1}(b)$, and its values at negative integers
are forced by the cocycle identity.  Every $b$ gives such a cocycle.
The equivalence and stabilizer formulas are exactly those in
Theorem~\ref{cat:J1}.
\end{proof}

\begin{proposition}[Refuted: slopes or weak admissibility do not imply uniqueness]\label{cat:J3}
The lead H1615-1 is false under the stated perfect-residue-field
hypothesis.  For $k=\mathbf F_p$, $L=K=\mathbf Q_p$, $G'=\mathbf G_m$,
and Witt Frobenius $\sigma=1$, the classes of
Corollary~\ref{cat:J2} are $\mathbf Q_p^\times$, whereas the Newton slope
is only $v_p:\mathbf Q_p^\times\to\mathbf Z$.
For every $a\in\mathbf Z_p^\times$, the rank-one filtered isocrystal
\[
 D_a=\mathbf Q_p e,\qquad \varphi(e)=a e,\qquad
 \operatorname{Fil}^iD_a=
 \begin{cases}D_a&i\leq0,\\0&i>0\end{cases}
\]
has slope zero and is weakly admissible.  Two such objects are
isomorphic if and only if their parameters are equal.  In particular
$a=1$ and $a=1+p$ give distinct classes with identical slopes and
identical filtration.
\end{proposition}
\begin{proof}
Every $\sigma$-coboundary is $h/\sigma(h)=1$, so the cyclic cohomology
quotient is the whole multiplicative group.  This quotient also occurs
with a nontrivial fixed subgroup: take any $\pi$ acting on $S$ with
$S^\pi=\mathbf Q_p$, extend its action to $\Gamma=\pi\times\mathbf Z$
by letting $\mathbf Z$ act trivially on $S$, put
$E=S^\times\rtimes\Gamma$, and fix
$\alpha_0(n)=(1,n)$.  The sections associated with $a$ are
$(n,m)\mapsto(a^m,(n,m))$.  All are continuous for the discrete topology
on $\mathbf Z$, they extend exactly the same $\alpha_0$, and
$H^\pi=\mathbf Q_p^\times$.  One may equally use a Frobenius lift on
$S$ in the reference section; its action on these invariants is still
identity.

For $D_a$, $t_N(D_a)=v_p(a)=0=t_H(D_a)$.  Its only proper subobject
is zero, so all weak-admissibility inequalities hold.  An isomorphism
$D_a\to D_b$ must send $e$ to $c e$ for some $c\in\mathbf Q_p^\times$.
Commutation with Frobenius requires $ac=bc$, whence $a=b$.
Thus neither equal slopes nor weak admissibility removes the ambiguity.
This is also a counterexample to conjugacy under the full group $H$,
by the last assertion of the first paragraph of
Theorem~\ref{cat:J1}'s proof.
\end{proof}

\begin{proposition}[Proved: the rank-one classification over a perfect residue field]\label{cat:J4}
Let $k$ be any perfect field of characteristic $p$, $L=W(k)[1/p]$, and
$\sigma$ Witt Frobenius.  Rank-one isocrystals over $L$ are classified by
\[
 L^\times/\{\sigma(c)/c:c\in L^\times\}.
\]
There is a split exact sequence of abelian groups
\[
 1\longrightarrow W(k)^\times/\{\sigma(u)/u:u\in W(k)^\times\}
 \longrightarrow H^1(\mathbf Z,L^\times)
 \xrightarrow{v_p}\mathbf Z\longrightarrow0,
\]
whose splitting sends $n$ to $[p^n]$.  If $k$ is algebraically closed,
the left term is trivial, so slopes classify rank-one isocrystals.
\end{proposition}
\begin{proof}
In a basis, Frobenius is $e\mapsto a e$ semilinearly, and replacing
the basis by $c e$ replaces $a$ by $a\sigma(c)/c$.  Each nonzero
$c\in L$ is uniquely $p^n u$ with $u\in W(k)^\times$, and
$\sigma(p)=p$.  The decomposition $L^\times=p^{\mathbf Z}\times
W(k)^\times$ proves exactness and splitting.

Suppose $k$ is algebraically closed and let $a\in W(k)^\times$.
Choose $c_1\in W(k)^\times$ such that
$\sigma(c_1)/c_1\equiv a\pmod p$; modulo $p$ this asks for a
nonzero root of $X^{p-1}=\overline a$.
Inductively assume the congruence holds modulo $p^n$, $n\geq1$.
Writing $c_{n+1}=c_n(1+p^n[x])$, the ratio changes modulo $p^{n+1}$
by the factor $1+p^n[x^p-x]$.
The required correction therefore reduces to one equation
$x^p-x=b$ in $k$, which has a root.  Choose such a correction at
each step.  The units $c_n$ form a $p$-adically Cauchy sequence;
its limit $c$ satisfies $\sigma(c)/c=a$.  This proves surjectivity
on units and the last assertion.  The argument works also for $p=2$.
\end{proof}

\subsection{The positive realization statement and its precise uniqueness}
\begin{theorem}[Known, with the hypotheses checked here]\label{cat:J5}
Let $\mathcal C_{\mathrm{BT}}$ be the smallest full Tannakian subcategory
of $\operatorname{Rep}_{\mathbf Q_p}(\pi_K)$ containing the rational
Tate modules of $p$-divisible groups over $\mathcal O_K$.
Then
\[
 D_{\mathrm{cris}}(V)=
 (B_{\mathrm{cris}}\otimes_{\mathbf Q_p}V)^{\pi_K}
\]
is an exact tensor functor on $\mathcal C_{\mathrm{BT}}$ to admissible
filtered $\varphi$-modules over $(L,K)$.  On Barsotti--Tate objects it
is their crystalline realization, with the covariant or contravariant
Dieudonne convention matched by dualizing as necessary.  Its de Rham
realization is $K\otimes_LD_{\mathrm{cris}}(V)$, and the associated
graded is the Hodge--Tate realization.
\end{theorem}
\begin{proof}[Derivation from opened theorems]
Fontaine~\cite[\S6.2]{catFontaine1982} proves that rational Tate
modules of $p$-divisible groups over $\mathcal O_K$ are crystalline
and identifies the filtered Dieudonne realization with the period
realization.  His theorem in \S5.2(i) proves that crystalline
representations are closed under the operations defining
$\mathcal C_{\mathrm{BT}}$; \S5.2(iii)--(v) gives the exact tensor
equivalence and the de Rham scalar-extension comparison.
The associated-graded comparison is the Hodge--Tate comparison of
\S3, in particular \S3.6--3.7 and \S3.10.  Thus all these assertions restrict to
$\mathcal C_{\mathrm{BT}}$.  We have used the actual BT-generated
subcategory; no assertion about all Hodge--Tate representations is
needed in this deduction.
\end{proof}

\begin{proposition}[Proved: uniqueness with the crystalline comparison]\label{cat:J6}
On $\mathcal C_{\mathrm{BT}}$, consider tensor functors $F$ to
admissible filtered $\varphi$-modules together with a specified tensor
natural isomorphism
$\eta:V_{\mathrm{cris}}\circ F\xrightarrow{\sim}
\operatorname{incl}_{\mathcal C_{\mathrm{BT}}}$.
The groupoid of such pairs is contractible: there is an isomorphism
to $(D_{\mathrm{cris}},\text{canonical comparison})$, unique among
isomorphisms compatible with $\eta$.  Weak admissibility may replace
admissibility in this statement.
\end{proposition}
\begin{proof}
Fontaine's equivalence in \S5.2 identifies
$D_{\mathrm{cris}}V_{\mathrm{cris}}F$ with $F$ by its counit.
Applying $D_{\mathrm{cris}}$ to $\eta$ and composing with the inverse
of this counit produces $F\simeq D_{\mathrm{cris}}|_{\mathcal C_{\mathrm{BT}}}$.
The construction is tensor and natural because the equivalence, counit,
and $\eta$ are.  Faithfulness of $V_{\mathrm{cris}}$ shows that at most
one such isomorphism induces the specified $\eta$.
Colmez--Fontaine~\cite[Theorem A; \S4.1, Theorem 4.3 and its
corollary]{catColmezFontaine2000} identifies weakly admissible
filtered $(\varphi,N)$-modules with admissible ones; setting $N=0$
gives the crystalline assertion.  This replacement preserves the
comparison datum, rather than deducing it from the Newton slopes.
\end{proof}

Kisin~\cite[Theorem (0.3), Corollary (2.2.6)]{catKisin2006} proves
that every crystalline representation of Hodge--Tate weights $0,1$
is the rational Tate module of a $p$-divisible group, with all primes
included at the isogeny level.  This identifies the weight-$[0,1]$
part of the crystalline category with the BT objects.  It does not
identify all Hodge--Tate representations of these weights with them.

The bare extension question on scan 024 asks for uniqueness from its
listed data and the restriction to BT objects.  Proposition~\ref{cat:J6}
proves uniqueness when the specified crystalline comparison to the given
Galois representations is part of those data.  It does not prove that
an unspecified extension has this comparison, or that an isomorphism
on the BT subcategory automatically respects every tensor relation in
the generated category.  Those extra implications, and the literal
larger category in the supplied scan-023 transcription, remain separate
gates.  The slope classification itself is completely refuted by
Proposition~\ref{cat:J3} and replaced by Theorem~\ref{cat:J1}.
